\begin{lemma}[A sub-window of the arc's own resolution, avoiding its truncated end edges, inherits $(\delta,\epsilon,n)$-ness]
  Let $E$ be a metric space, $\gamma \in \Theta(E)$ (\noderef{Curve}) a closed curve
  (\noderef{IsClosedCurve}), $\delta,\epsilon \in \mathbb{R}$, $n,k \in \mathbb{N}$, and $s \colon
  \mathbb{N} \to \mathbb{R}$ a geodesic resolution of $\gamma$ into $k$ edges
  (\noderef{IsGeodesicResolution}) such that, for this particular $s$, both counting clauses of
  $\mathrm{IsDeltaEpsN}$ (\noderef{IsDeltaEpsN}) hold:
  \[
    \text{(linear)}\quad \forall\, a,b,\ 0 \le a \le b \le \length(\gamma),\ b-a \le
    2\epsilon^{-1}\delta \implies \#\{j\in\set{1,\dots,k} : a\le s(j-1),\ s(j)\le b,\ s(j)-s(j-1)<\delta\}
    \le n
    ,
  \]
  \[
    \text{(wrap)}\quad \forall\, a,b,\ 0 \le b < a \le \length(\gamma),\
    (\length(\gamma)-a)+b \le 2\epsilon^{-1}\delta \implies
    \#\{j\in\set{1,\dots,k} : (a\le s(j-1) \vee s(j)\le b),\ s(j)-s(j-1)<\delta\} \le n
    .
  \]
  Let $t,t' \in [0,\length(\gamma)]$, and let $(k_A,s_A)$ be any geodesic resolution of the
  circular arc $A := \gamma|_{[t,t']}$ (\noderef{curveArc}) with $k_A \ge 1$ carrying
  \noderef{ArcGeodesicResolution}'s correspondence: in the ordinary case $t \le t'$, there is
  $j_0$ with $s_A(i) = s(j_0+(i-1))-t$ for every $i \in \set{1,\dots,k_A-1}$ (with
  $j_0+(i-1)\le k$); in the wrap case $t' < t$, there are $j_0,m$ with $s_A(i) = s(j_0+(i-1))-t$
  for $1 \le i \le m$ and $s_A(i) = (\length(\gamma)-t)+s(i-m)$ for $m < i \le k_A-1$ (with the
  stated index bounds), \emph{together with the fact that the two halves of this correspondence
  meet at a genuine edge of $\gamma$}: $s(j_0+(m-1)) = \length(\gamma)$ whenever $m \ge 1$. (This
  extra fact is now itself part of \noderef{ArcGeodesicResolution}'s stated wrap-case conclusion --
  see the remark below -- so a call site instantiating $(k_A,s_A,j_0,m)$ from that lemma discharges
  it directly from the cited conclusion, with no re-derivation; it is still recorded here as an
  explicit hypothesis of the present lemma's own signature, since this lemma is stated for an
  arbitrary resolution carrying the correspondence shape, not tied to
  \noderef{ArcGeodesicResolution}'s specific construction.)

  Suppose the sub-window of $A$ strictly between its own breakpoints $u := s_A(1)$ and $u' :=
  s_A(k_A-1)$, i.e.\ $A|_{[u,u']} := \mathrm{curveRestrict}(A,u,u')$ (\noderef{curveRestrict}), is
  \emph{not} a closed curve (\noderef{IsClosedCurve}). Then $A|_{[u,u']}$ is a
  $(\delta,\epsilon,n)$-curve (\noderef{IsDeltaEpsN}).

  \paragraph{Why the sub-window, not the whole arc.} paper.tex 899--901 asserts, in the setting of
  \Cref{typeI-cut-point-exists}, that ``$\gamma|_{[u,u']}$'' -- the arc with its up to two partial
  end edges dropped -- inherits $(\delta,\epsilon,n)$-ness from $\gamma$. The whole arc $A$ does
  \emph{not} inherit it: its first and last edges are the (possibly) truncated partial edges
  supplied by \noderef{ArcGeodesicResolution}, and a truncated edge can be short for every
  $\delta>0$ regardless of $n$, exactly as in the counterexample recorded on
  \noderef{RestrictDeltaEpsNAtBreakpoints}. Excluding exactly those two edges -- the window between
  the arc's own interior breakpoints $s_A(1)$ and $s_A(k_A-1)$ -- is what makes the inheritance
  true, and this is the shape \Cref{cut-type-I} Branch~C consumes directly.

  \paragraph{Why this is the sole consumer of $\mathrm{IsDeltaEpsN}$'s wrap clause.} As of this
  cycle, the (wrap) clause recorded above has no consumer anywhere in the tablet and is discharged
  vacuously at every other call site. It is needed here because, in the wrap case $t'<t$, a window
  of the sub-window can straddle the point where $A$'s tail piece $\gamma|_{[t,\length(\gamma)]}$
  meets its head piece $\gamma|_{[0,t']}$ -- a circular window of $\gamma$ straddling the basepoint
  $\gamma(0)=\gamma(\length(\gamma))$, exactly the situation the (wrap) clause was written for.

  \paragraph{On the hypothesis $s(j_0+(m-1))=\length(\gamma)$.} This is true of the specific
  resolution \noderef{ArcGeodesicResolution} constructs (its tail piece $\gamma|_{[t,\length(\gamma)]}$
  has length exactly $\length(\gamma)-t$, so its own resolution's top breakpoint -- glued to $s_A$ at
  index $m$ -- necessarily lands there), and, as of this cycle, \noderef{ArcGeodesicResolution}'s
  stated wrap-case conclusion exposes it directly as its third conjunct. Without this fact, the edge
  $[s_A(m-1),s_A(m)]$ and the edge $[s_A(m),s_A(m+1)]$ meeting at the split could fail to correspond
  to genuine, single edges of $\gamma$, and the counting argument below would not close; that is why
  it is needed here at all. A call site instantiating $(k_A,s_A,j_0,m)$ from
  \noderef{ArcGeodesicResolution} therefore obtains this hypothesis directly from that lemma's own
  stated conclusion, with no re-derivation -- it is no longer a gap awaiting a future
  strengthening. It is kept as an explicit hypothesis of the present lemma's signature because the
  present lemma is stated for an arbitrary resolution carrying the correspondence shape, and does
  not itself assume that resolution came from \noderef{ArcGeodesicResolution}.

  \paragraph{The general (non-breakpoint) restriction of $A$ does not inherit this.} Exactly as for
  \noderef{RestrictDeltaEpsNAtBreakpoints}, restricting $A$ to an arbitrary sub-window (not a pair
  of $s_A$'s own breakpoints) can truncate an edge of $s_A$ mid-edge; and restricting $A$ to all of
  $[0,\length(A)]$ (i.e.\ not dropping the two end edges) is exactly the scenario
  \noderef{RestrictDeltaEpsNAtBreakpoints} shows can fail, applied to a truncated end edge of $A$
  itself.
\end{lemma}
\begin{proof}
  Throughout write $A := \gamma|_{[t,t']}$ (\noderef{curveArc}), $u := s_A(1)$, $u' := s_A(k_A-1)$
  and $R := A|_{[u,u']} = \mathrm{curveRestrict}(A,u,u')$ (\noderef{curveRestrict}). By
  \noderef{IsGeodesicResolution} applied to $(k,s)$: $s$ is monotone on $\set{0,\dots,k}$,
  $s(0)=0$ and $s(k)=\length(\gamma)$; hence
  \begin{equation}\label{ADEN-srange}
    0 \le s(j) \le \length(\gamma) \qquad \text{for every } j \in \set{0,\dots,k},
  \end{equation}
  by monotonicity from $s(0)=0$ resp.\ to $s(k)=\length(\gamma)$. Likewise $s_A$ is monotone on
  $\set{0,\dots,k_A}$ with $s_A(0)=0$ and $s_A(k_A)=\length(A)$.

  \textbf{Step 0: $k_A=1$ is impossible.} If $k_A=1$ then $s_A(k_A-1)=s_A(0)=0$, and the standing
  hypotheses $0\le s_A(1)$ and $s_A(1)\le s_A(k_A-1)=0$ force $s_A(1)=0$. Hence
  $\length(R)=u'-u=0-0=0$ (\noderef{curveRestrict}'s $\length$ field is definitionally $b-a$), so
  $R.\mathrm{toFun}(\length(R)) = R.\mathrm{toFun}(0)$ literally, i.e.\ $R$ \emph{is} a closed
  curve (\noderef{IsClosedCurve}), contradicting the standing hypothesis that it is not. So
  $k_A\ge2$ from here on.

  \textbf{Step 1: the resolution witness.} Since $k_A\ge2$ we have $1\le k_A-1\le k_A$, and
  $0\le u$, $u\le u'$, $u'\le\length(A)$ are the standing hypotheses; so
  \noderef{RestrictPiecewiseGeodesic}, applied to $(k_A,s_A)$ at $p=1$, $q=k_A-1$, says that
  \[
    s'(j) := s_A(1+j) - u
  \]
  is a geodesic resolution of $R$ into $(k_A-1)-1$ edges, i.e.\
  $\mathrm{IsGeodesicResolution}(R,\,k_A-2,\,s')$. We use $(k_A-2,s')$ as the witness for
  $\mathrm{IsDeltaEpsN}(R,\delta,\epsilon,n)$ (\noderef{IsDeltaEpsN}) and verify its two remaining
  conjuncts.

  \textbf{Step 2: the closed-curve (wrap) clause of $\mathrm{IsDeltaEpsN}$ is vacuous.} That clause
  is guarded by ``$R$ is a closed curve'', and $R$ is hypothesised not to be
  (\noderef{IsClosedCurve}). So there is nothing to prove.

  \textbf{Step 3: reduction of the linear clause to a windowed count on $s_A$.} Fix $a,b$ with
  $0\le a\le b\le\length(R)=u'-u$ and $b-a\le2\epsilon^{-1}\delta$. We must bound
  \[
    S := \#\set{ j\in\set{1,\dots,k_A-2} \colon a\le s'(j-1),\ s'(j)\le b,\ s'(j)-s'(j-1)<\delta }
    .
  \]
  For $j\ge1$ we have $1+(j-1)=j$, so $s'(j-1)=s_A(j)-u$, while $s'(j)=s_A(1+j)-u=s_A(j+1)-u$.
  Thus the three conditions on such a $j$ read, after adding $u$ throughout and writing $i:=j+1$
  (so that $j = i-1$ and $2\le i\le k_A-1$):
  \[
    a+u \le s_A(i-1), \qquad s_A(i) \le b+u, \qquad s_A(i)-s_A(i-1) < \delta .
  \]
  The map $j\mapsto j+1$ is injective, so
  \begin{equation}\label{ADEN-reindex}
    S \;\le\; N := \#\set{ i\in\set{2,\dots,k_A-1} \colon a' \le s_A(i-1),\ s_A(i)\le b',\
      s_A(i)-s_A(i-1)<\delta },
  \end{equation}
  where from now on $a' := a+u$ and $b' := b+u$. It remains to prove $N\le n$.

  \textbf{Step 4: the clamped-linear bound.} By \eqref{ADEN-srange}, the monotonicity of $s$ on
  $\set{0,\dots,k}$, and the hypothesis (linear), \noderef{ClampedLinearShortEdgeCount} applies
  with $L:=\length(\gamma)$ and yields: for \emph{every} pair $x,y\in\mathbb{R}$ with
  $y-x\le2\epsilon^{-1}\delta$,
  \begin{equation}\label{ADEN-clamp}
    \#\set{ j\in\set{1,\dots,k} \colon x\le s(j-1),\ s(j)\le y,\ s(j)-s(j-1)<\delta } \le n
    .
  \end{equation}
  We use this three times below; each time the pair $(x,y)$ is $(a'+\theta,\,b'+\theta)$ for a
  single real shift $\theta$, whose gap is $y-x = b'-a' = b-a \le 2\epsilon^{-1}\delta$ by the
  standing hypothesis on $(a,b)$, so \eqref{ADEN-clamp} is applicable each time.

  \textbf{Step 5: the ordinary case $t\le t'$.} Take $j_0$ from the hypothesis: for
  $1\le i\le k_A-1$, $j_0+(i-1)\le k$ and $s_A(i)=s(j_0+(i-1))-t$. Since $i\ge1$ gives
  $i+j_0-1=j_0+(i-1)$, this is exactly the affine-correspondence hypothesis of
  \noderef{ShiftedBreakpointShortEdgeCount} with parameters
  \[
    (p,q,c,d,\sigma,a',b') = (2,\ k_A-1,\ j_0,\ 1,\ -t,\ a',\ b'),
  \]
  whose side conditions $1\le p$ and $d+1\le p+c$ read $1\le2$ and $1+1\le2+j_0$, both true, and
  whose correspondence range $\set{p-1,\dots,q}=\set{1,\dots,k_A-1}$ is exactly the range where
  the hypothesis supplies the formula. Its remaining hypothesis (window count) is
  \eqref{ADEN-clamp} at $(x,y)=(a'-(-t),\,b'-(-t))=(a'+t,\,b'+t)$. Its conclusion is precisely
  $N\le n$.

  \textbf{Step 6: the wrap case $t'<t$.} Take $j_0,m$ from the hypothesis, with the tail formula
  $s_A(i)=s(j_0+(i-1))-t$ for $1\le i\le m$, the head formula
  $s_A(i)=(\length(\gamma)-t)+s(i-m)$ for $m<i\le k_A-1$, and the splice identity
  $s(j_0+(m-1))=\length(\gamma)$ whenever $m\ge1$.

  First a length fact. Unfolding \noderef{curveArc}'s wrap branch together with
  \noderef{curveWrapRestrict}, \noderef{curveConcat} and \noderef{curveRestrict}'s length fields
  (all definitional), $\length(A)=(\length(\gamma)-t)+(t'-0)$, and $t'<t$ gives
  \begin{equation}\label{ADEN-margin}
    \length(A) < \length(\gamma)
    .
  \end{equation}
  Since $0\le a\le b\le u'-u$ and $0\le u\le u'\le\length(A)$, we get
  $b-a\le b\le u'-u\le\length(A)$, so by \eqref{ADEN-margin}
  \begin{equation}\label{ADEN-gapsmall}
    b-a < \length(\gamma)
    .
  \end{equation}

  \emph{The head bound.} Apply \noderef{ShiftedBreakpointShortEdgeCount} with
  \[
    (p,q,c,d,\sigma,a',b') = \bigl(\max(2,m+1),\ k_A-1,\ 0,\ m,\ \length(\gamma)-t,\ a',\ b'\bigr).
  \]
  Side conditions: $1\le\max(2,m+1)$ and $m+1\le\max(2,m+1)+0$, both true. Correspondence range:
  $\set{\max(2,m+1)-1,\dots,k_A-1}$. Let $i$ be in that range. If $i>m$, the head formula gives
  $i-m\le k$ and $s_A(i)=s(i-m)+(\length(\gamma)-t)$, and $i+0-m=i-m$, as required. If instead
  $i\le m$, then $\max(2,m+1)-1\le i\le m$ forces $i=m$ and $m\ge1$ (because
  $\max(2,m+1)-1\ge m$ and $\max(2,m+1)-1\ge1$); then $i+0-m=0$ and, by the tail formula at $i=m$
  (legitimate as $1\le m\le m$) together with the splice identity (legitimate as $m\ge1$) and
  $s(0)=0$,
  \[
    s_A(m) = s\bigl(j_0+(m-1)\bigr)-t = \length(\gamma)-t = s(0)+\bigl(\length(\gamma)-t\bigr),
  \]
  again as required, with index bound $0\le k$. The window-count hypothesis is \eqref{ADEN-clamp}
  at $(x,y)=\bigl(a'-(\length(\gamma)-t),\ b'-(\length(\gamma)-t)\bigr)$. We obtain
  \begin{equation}\label{ADEN-head}
    N_H := \#\set{ i\in\set{\max(2,m+1),\dots,k_A-1} \colon a'\le s_A(i-1),\ s_A(i)\le b',\
      s_A(i)-s_A(i-1)<\delta } \le n .
  \end{equation}

  \emph{The tail bound.} Apply \noderef{ShiftedBreakpointShortEdgeCount} with
  $(p,q,c,d,\sigma,a',b')=(2,\ m,\ j_0,\ 1,\ -t,\ a',\ b')$. Side conditions $1\le2$ and
  $1+1\le2+j_0$ hold; the correspondence range is $\set{1,\dots,m}$, exactly where the tail
  formula is supplied (and $i+j_0-1=j_0+(i-1)$ for $i\ge1$); the window count is
  \eqref{ADEN-clamp} at $(x,y)=(a'+t,\,b'+t)$. We obtain
  \begin{equation}\label{ADEN-tail}
    N_T := \#\set{ i\in\set{2,\dots,m} \colon a'\le s_A(i-1),\ s_A(i)\le b',\
      s_A(i)-s_A(i-1)<\delta } \le n .
  \end{equation}

  \emph{Covering.} Write $F$, $F_T$, $F_H$ for the three index sets counted by $N$, $N_T$, $N_H$
  respectively (all three carry the same three defining conditions; only the ambient index range
  differs). Then $F\subseteq F_T\cup F_H$: if $i\in F$ then $2\le i\le k_A-1$, and either $i\le m$,
  putting $i\in F_T$, or $i\ge m+1$, which with $i\ge2$ gives $i\ge\max(2,m+1)$ and puts
  $i\in F_H$.

  \emph{Sub-case $F_T=\emptyset$.} Then $F\subseteq F_H$, so $N\le N_H\le n$ by \eqref{ADEN-head}.

  \emph{Sub-case $F_H=\emptyset$.} Then $F\subseteq F_T$, so $N\le N_T\le n$ by \eqref{ADEN-tail}.

  \emph{Sub-case $F_T\ne\emptyset$ and $F_H\ne\emptyset$.} Fix $i_T\in F_T$ and $i_H\in F_H$; so
  $2\le i_T\le m$ (whence $m\ge2\ge1$) and $\max(2,m+1)\le i_H\le k_A-1$ (whence $m+1\le k_A-1$,
  i.e.\ $m+2\le k_A$, and $i_H>m$). We apply \noderef{SplicedBreakpointShortEdgeCount} with
  \[
    (K,c,m,\sigma,\tau,a',b') = \bigl(k_A,\ j_0,\ m,\ -t,\ \length(\gamma)-t,\ a',\ b'\bigr)
  \]
  and check each of its hypotheses.
  \begin{itemize}
    \item $s$ monotone on $\set{0,\dots,k}$ and $s(0)=0$: \noderef{IsGeodesicResolution} for
      $(k,s)$.
    \item $1\le m$ and $m+2\le k_A$: just established.
    \item Tail correspondence, for $1\le i\le m$: $j_0+(i-1)\le k$ and
      $s_A(i)=s(j_0+(i-1))-t=s(j_0+(i-1))+\sigma$ with $\sigma=-t$. This is the hypothesis
      verbatim.
    \item Head correspondence, for $m<i\le k_A-1$: $i-m\le k$ and
      $s_A(i)=(\length(\gamma)-t)+s(i-m)=s(i-m)+\tau$ with $\tau=\length(\gamma)-t$. Again the
      hypothesis verbatim.
    \item Splice: $s(j_0+(m-1))+\sigma = \length(\gamma)-t = \tau$, using the splice identity
      (available since $m\ge1$).
    \item Margin: $b'-\tau < a'-\sigma$ reads $(b+u)-(\length(\gamma)-t) < (a+u)+t$, i.e.\
      $b-a<\length(\gamma)$, which is \eqref{ADEN-gapsmall}.
    \item Circular window count: we must bound
      $\#\set{j\in\set{1,\dots,k} : (\alpha\le s(j-1)\vee s(j)\le\beta),\ s(j)-s(j-1)<\delta}$
      by $n$, where $\alpha := a'-\sigma = a'+t$ and $\beta := b'-\tau = b'-(\length(\gamma)-t)$.
      This is the hypothesis (wrap) at the pair $(\alpha,\beta)$, whose four side conditions we
      verify.
      \begin{itemize}
        \item $0\le\beta$: since $i_H>m$ and $i_H\le k_A-1$, the head formula gives $i_H-m\le k$
          and $s(i_H-m)=s_A(i_H)-(\length(\gamma)-t)$. As $i_H\in F_H$ we have $s_A(i_H)\le b'$,
          so $s(i_H-m)\le b'-(\length(\gamma)-t)=\beta$; and $s(i_H-m)\ge0$ by
          \eqref{ADEN-srange}. Hence $0\le\beta$.
        \item $\alpha\le\length(\gamma)$: since $2\le i_T\le m$ we have $1\le i_T-1\le m$, so the
          tail formula applies at $i_T-1$ and gives $j_0+(i_T-2)\le k$ and
          $s(j_0+(i_T-2))=s_A(i_T-1)+t$. As $i_T\in F_T$ we have $a'\le s_A(i_T-1)$, so
          $\alpha=a'+t\le s_A(i_T-1)+t=s(j_0+(i_T-2))$; and $s(j_0+(i_T-2))\le\length(\gamma)$ by
          \eqref{ADEN-srange}. Hence $\alpha\le\length(\gamma)$.
        \item $\beta<\alpha$: equivalent to $b-a<\length(\gamma)$, which is
          \eqref{ADEN-gapsmall}.
        \item $(\length(\gamma)-\alpha)+\beta\le2\epsilon^{-1}\delta$: indeed
          $(\length(\gamma)-a-u-t)+(b+u-\length(\gamma)+t) = b-a \le 2\epsilon^{-1}\delta$.
      \end{itemize}
  \end{itemize}
  The conclusion of \noderef{SplicedBreakpointShortEdgeCount} is exactly $N\le n$.

  In every case $N\le n$, so $S\le n$ by \eqref{ADEN-reindex}. This is the linear counting clause
  for the witness $(k_A-2,s')$, and with Steps~1 and 2 it completes the verification of
  $\mathrm{IsDeltaEpsN}(R,\delta,\epsilon,n)$.
\end{proof}
